% =====================================================================
% Test de non-regression : symbole de fin de boite sous beamer.
%
% Le symbole (□ exemple, ■ preuve) ne doit apparaitre que sur la DERNIERE
% diapositive de la boite, quelle que soit la facon de la decouper. Verifie par
% check-qed-overlays.sh (une page = un overlay).
%   frame 1 : \only<1>..<3>   -> symbole sur l'overlay 3 seulement
%   frame 2 : \onslide<k->   + \qedhere -> symbole inline sur l'overlay 4 seulement
%   frame 3 : preuve en \only<1>/<2>     -> overlay 2 seulement
%   frame 4 : \pause seul                -> overlay 2 seulement
%   frame 5 : sans overlay               -> symbole present
% =====================================================================
\documentclass[9pt,t]{beamer}
\usepackage[lang=fr, theme=ocots, math=analysis]{ocots}

\begin{document}

\begin{frame}{Exemple découpé en only}
\begin{example*}[title={Découpé en only}]
  \only<1>{Premier pas.}
  \only<2>{Deuxième pas.}
  \only<3>{Troisième pas.}
\end{example*}
\end{frame}

\begin{frame}{Onslide et qedhere}
\begin{example*}[title={Align}]
  \begin{align*}
    a &= 1 \\
    \onslide<2->{b &= 2 \\}
    \onslide<3->{c &= 3 \\}
    \onslide<4->{d &= 4 \qedhere}
  \end{align*}
\end{example*}
\end{frame}

\begin{frame}{Preuve en only}
\begin{proof}
  \only<1>{Début de la preuve.}
  \only<2>{Fin de la preuve.}
\end{proof}
\end{frame}

\begin{frame}{Pause seul}
\begin{example}
  Un premier point.
  \pause
  Un second point.
\end{example}
\end{frame}

\begin{frame}{Sans overlay}
\begin{example}
  Une boîte ordinaire.
\end{example}
\begin{proof}
  Une preuve ordinaire.
\end{proof}
\end{frame}

\end{document}
